% English translation, made from our own transcription (original.tex), never from any existing
% translation — and there was none to be tempted by: Hirzebruch, looking for a Noether publication
% on this subject, could point only at this abstract, and noted it was not even reprinted in her
% Gesammelte Abhandlungen. The \origpage marker, the inline normal basis and both bare italic e's
% are reproduced unchanged; only the prose is translated.
%
% THE ONE BRACKETED ADDITION. "als direkte Summe größter zyklischer" is elliptical in the print —
% the head noun is simply not set, which German carries comfortably and English does not. It is
% supplied in square brackets, "largest cyclic [groups]", following the translate prompt's rule
% that a genuinely necessary gloss is bracketed and never parenthesized, parentheses being the
% author's own device for asides. Nothing else is added, and nothing is removed.
%
% PERIOD TERMS KEPT LITERAL rather than modernized, because the whole point of the abstract is
% which theorem is prior to which, so the reader needs to meet Noether's own words for the
% objects: "residue class system" for Restklassensystem, not "quotient group"; "modules of
% integral linear forms" for Moduln aus ganzzahligen Linearformen, not "lattices"; "normal basis"
% for Normalbasis, not "Smith normal form". See glossary.yaml.
%
% THE TWO EM-DASHES are Noether's own parenthetical dashes inside her final sentence, not the
% report's speaker separators, and are kept. The abbreviation "z.\ B." becomes "e.\ g." with the
% same control space, matching how the transcription sets it.
\documentclass{article}
\usepackage{readmasters}
\begin{document}

\origpage{104}

\section*{Derivation of elementary divisor theory from group theory.}

\subsection*{E. Noether.}

As is well known, elementary divisor theory gives, for modules of integral linear forms, a normal basis of the form $(e_1 y_1, e_2 y_2, \ldots, e_r y_r)$, where each $e$ is divisible by the one following; the $e$ are thereby uniquely determined up to sign. Since every Abelian group with finitely many generators is isomorphic to the residue class system modulo such a module, the decomposition theorem for these groups as a direct sum of largest cyclic [groups] is thereby proved as well. Conversely, the decomposition theorem is now obtained directly and purely group-theoretically, in generalization of the proof customary for finite groups, and from it elementary divisor theory is derived by passing from the residue class system to the module itself. The group theorem thus proves to be the simpler theorem; in the applications of the group theorem --- e.\ g. Betti and torsion numbers in topology --- going back to elementary divisor theory is accordingly not required.

\end{document}
